% Generated mechanically from the frozen guide.tex target.
% Do not edit this generated file; edit the bound locale source instead.

% OK, start here.
%

\setcounter{chapter}{111}
\chapter{文献指南}


\phantomsection
\label{guide-section-phantom}


\section{简短导论文章}
\label{guide-section-short-introductions}


\begin{itemize}
\item Barbara Fantechi: \emph{Stacks for Everybody} \cite{fantechi_stacks}
\item Dan Edidin: \emph{What is a stack?} \cite{edidin_whatis}
\item Dan Edidin: \emph{Notes on the construction of the moduli space of
curves} \cite{edidin_notes}
\item Angelo Vistoli: \emph{Intersection theory on algebraic
stacks and on their moduli spaces}，尤其是附录。
\cite{vistoli_intersection}
\end{itemize}


\section{经典参考文献}
\label{guide-section-classics}

\begin{itemize}
\item Mumford: \emph{Picard groups of moduli problems}
\cite{mumford_picard}
\begin{quote}
本文从未使用 ''stack'' 一词，但这一概念隐含其中；作者计算了椭圆曲线模叠的 Picard 群。
\end{quote}
\item Deligne, Mumford: \emph{The irreducibility of the space of curves of
given genus} \cite{DM}
\begin{quote}
这篇有影响力的论文引入了 ''algebraic stacks''；如今通常称为 Deligne--Mumford 叠（具有可表示对角且允许由概形给出的 \'etale 展示的叠）。其中有许多 \emph{未给出证明} 的基础性结果。论文利用叠给出了亏格为 $g$ 的曲线模空间不可约性的两个证明。
\end{quote}
\item Artin: \emph{Versal deformations and algebraic stacks}
\cite{ArtinVersal}
\begin{quote}
本文引入了推广 Deligne--Mumford 叠的 ''algebraic stacks''，现在通常称为 \emph{Artin stacks}；其对角可表示并允许由概形给出的光滑展示。本文给出称为 Artin 判据的形变论判据，可在不显式构造展示的情况下证明给定模叠是 Artin 叠。
\end{quote}
\end{itemize}




\section{书籍与在线讲义}
\label{guide-section-books}


\begin{itemize}
\item Laumon, Moret-Bailly: \emph{Champs Alg\'ebriques} \cite{LM-B}
\begin{quote}
这是目前关于叠且包含许多基础结果的最详尽参考书。它假定读者熟悉代数空间，并频繁引用 Knutson 的书 \cite{Kn}。第 12 章关于代数叠 lisse-\'etale site 的函子性存在一个错误。无需担心，因为 Martin Olsson 已修补该错误（见 \cite{olsson_sheaves}），其余各章的结果（可能需要略作修改）都是正确的。
\end{quote}
\item The Stacks Project Authors: \emph{Stacks Project}
\cite{stacks-project}.
\begin{quote}
你正在阅读它！
\end{quote}
\item Anton Geraschenko:
\emph{Lecture notes for Martin Olsson's class on stacks} \cite{olsson_stacks}
\begin{quote}
本课程先系统发展代数空间理论，再引入代数叠（第 27 讲才首次定义！）。除基本性质外，课程涵盖 Deligne--Mumford 条件与具有非分歧对角之间的等价性、Artin 叠上的 lisse-\'etale site、拟凝聚层理论、Keel--Mori 定理、上同调下降以及 gerbe（及其与 Brauer 群的关系）。另有一些练习。
\end{quote}
\item
Behrend, Conrad, Edidin, Fantechi, Fulton, G\"ottsche, and Kresch:
\emph{Algebraic stacks}, online notes for a book being currently written
\cite{stacks_book}
\begin{quote}
本书旨在友好地介绍叠，重点放在例子与应用上，不要求读者具有高深背景。不同于 \cite{LM-B}，它不假定读者已经掌握代数空间理论；相反，先引入 Deligne--Mumford 叠，并将代数空间作为特殊情形，目标之一是发展足够理论以证明 \cite{DM} 中的断言。第二部分将发展 Artin 叠的一般理论。目前 Kresch 网站上只有书的一小部分可用。
\end{quote}
\item Olsson, Martin: \emph{Algebraic spaces and stacks}, \cite{olsson_book}
\begin{quote}
这是代数空间与代数叠的高度推荐入门读物，起点是已经掌握 Hartshorne 代数几何的读者。
\end{quote}
\end{itemize}


\section{关于叠基础的相关参考文献}
\label{guide-section-related}


\begin{itemize}
\item Vistoli:
\emph{Notes on Grothendieck topologies, fibered categories and descent theory}
\cite{vistoli_fga}
\begin{quote}
包含关于纤维范畴、叠以及 fpqc 拓扑中的下降理论的有用事实，并给出严谨证明。
\end{quote}
\item Knutson: \emph{Algebraic Spaces} \cite{Kn}
\begin{quote}
本书源自作者在 Michael Artin 指导下的博士论文，包含代数空间理论的基础。书 \cite{LM-B} 经常引用它。另见 Artin 关于代数空间的论文：\cite{Artin-Algebraic-Approximation}、
\cite{ArtinI}、\cite{Artin-Implicit-Function}、
\cite{ArtinII}、\cite{Artin-Construction-Techniques}、
\cite{Artin-Algebraic-Spaces}、\cite{Artin-Theorem-Representability} 以及
\cite{ArtinVersal}
\end{quote}
\item Grothendieck et al, \emph{Th\'eorie des Topos et Cohomologie \'Etale des
Sch\'emas I, II, III}，亦称 SGA4 \cite{SGA4}
\begin{quote}
第 1 卷包含关于宇宙、site 与纤维范畴的许多一般事实。''champ'' 一词（法语中表示 ''stack''）出现在 Deligne 的 Expos\'e XVIII。
\end{quote}
\item Jean Giraud: \emph{Cohomologie non ab\'elienne} \cite{giraud}
\begin{quote}
本书讨论一般 site 上的纤维范畴、叠、torsor 与 gerbe，但不讨论代数叠。例如，若 $G$ 是 $X$ 上的阿贝尔群层，则如同 $H^1(X, G)$ 可与 $G$-torsor 识别一样，$H^2(X, G)$ 可与适当定义的 $G$-gerbe 集合识别。当 $G$ 非阿贝尔时，$H^2(X, G)$ 定义为 $G$-gerbe 的集合。
\end{quote}
\item Kelly and Street: \emph{Review of the elements of 2-categories}
\cite{kelly-street}
\begin{quote}
叠的范畴形成一个 2-范畴，虽然后者是 2-范畴的一种简单类型，其中 2-态射都是可逆的。这是关于一般 2-范畴的参考文献。作者没有使用过它，因而无法评价其有用程度；另请注意 \cite{stacks-project} 包含一些关于 2-范畴的基础内容。
\end{quote}
\end{itemize}




\section{文献中的论文}
\label{guide-section-papers}

\noindent
下面列出包含叠与代数空间基本结果的研究论文。摘要的目的仅是指出对叠理论有贡献的结果；在许多情况下，这些结果只是论文主要目标的附带结果。我们将论文分为若干类别，有些论文会属于多个类别。


\subsection{形变理论与代数叠}
\label{guide-subsection-deformation-theory}

\noindent
Artin 的前三篇论文虽不涉及叠，却包含强有力的结果；前两篇对
\cite{ArtinVersal} 至关重要。
\begin{itemize}
\item Artin: \emph{Algebraic approximation of structures over
complete local rings} \cite{Artin-Algebraic-Approximation}
\begin{quote}
论文证明，在适当的温和假设下，任意有效形式形变都可以近似：若
$F: (\Sch/S) \to (\textit{Sets})$ 是一个逆变函子，
在有限型域或优良 DVR 上局部有限表示（基概形为 $S$），且 $s \in S$，
$\hat{\xi} \in F(\hat{\mathcal{O}}_{S, s})$ 是有效形式形变，
则对任意 $n > 0$，存在残余平凡的 \'etale 邻域
$(S', s') \to (S, s)$ 以及 $\xi' \in F(S')$，使得 $\xi'$ 与
$\hat{\xi}$ 一致到 $n$ 阶（即在
$F(\mathcal{O}_{S, s} / \mathfrak m^n)$ 中具有相同限制）。
\end{quote}
\item
Artin: \emph{Algebraization of formal moduli I} \cite{ArtinI}
\begin{quote}
论文证明，在温和假设下任意有效形式 versal 形变都可代数化。设
$F: (\Sch/S) \to (\textit{Sets})$ 为逆变函子，在有限型域或优良 DVR
上局部有限表示（基概形为 $S$），$s \in S$ 是局部闭点，$\hat A$ 是完备 Noether 局部
$\mathcal{O}_S$-代数，剩余域 $k'$ 是 $k(s)$ 的有限扩张，
并且 $\hat{\xi} \in F(\hat A)$ 是元素 $\xi_0 \in F(k')$ 的有效形式
versal 形变。则存在有限型 $S$-概形 $X$、闭点 $x \in X$（剩余域
$k(x) = k'$）以及 $\xi \in F(X)$，存在同构
$\hat{\mathcal{O}}_{X, x} \cong \hat{A}$，并且在每个
$F(\hat A / \mathfrak m^n)$ 中识别 $\xi$ 与 $\hat{\xi}$ 的限制。
若 $\hat{\xi}$ 是普遍形变，则代数化是唯一的。文中还给出了 Hilbert
概形与 Picard 概形可表示性的应用。
\end{quote}
\item Artin: \emph{Algebraization of formal moduli. II}
\cite{ArtinII}
\begin{quote}
粗略地说，论文证明：若能在 $Y' \subset X'$ 附近形式局部地收缩闭子集
$Y'$，则存在全局态射 $X' \to X$ 收缩 $Y$，其中 $X$ 是代数空间。
\end{quote}
\item
Artin: \emph{Versal deformations and algebraic stacks} \cite{ArtinVersal}
\begin{quote}
这篇重要论文建立在
\cite{Artin-Algebraic-Approximation} 与 \cite{ArtinI} 的工作之上。
论文引入 Artin 判据：通过检验形变论性质即可证明叠的代数性。更准确地说
（但仍非完全严格地说），Artin 在点 $x \in \mathcal{X}(k)$ 附近为极限保持叠
$\mathcal{X}$ 构造如下展示：假设叠 $\mathcal{X}$ 满足 Schlessinger
判据（\cite{Sch}），则存在 $x$ 的形式 versal 形变
$\hat{\xi} \in \lim \mathcal{X}(\hat A / \mathfrak m^n)$。若形式形变有效
（即 $\mathcal{X}(\hat{A}) \to \lim \mathcal{X}(\hat A / \mathfrak m^n)$
为双射），则得到有效形式 versal 形变 $\xi \in \mathcal{X}(\hat A)$。
利用 \cite{ArtinI} 的结果，可取有限型概形 $U$ 及元素
$\xi_U: U \to \mathcal{X}$，它在 $x$ 上方的点 $u \in U$ 处形式 versal。
若进一步假设 $\mathcal{X}$ 具有满足某些条件的形变与障碍理论（即与
\'etale 局部化、完备化兼容且满足可构造性条件），则第 4 节证明形式 versal
性是开条件；因此缩小 $U$ 后，$U \to \mathcal{X}$ 光滑。
Artin 还证明，任何具有概形 fppf 展示的叠都有概形光滑展示；特别地，
可用平坦、分离、有限表示的群概形构造商叠。
\end{quote}
\item Conrad, de Jong: \emph{Approximation of Versal Deformations}
\cite{conrad-dejong}
\begin{quote}
本文通过应用 Popescu 的强结果研究 Artin 的代数化结果：若 $A$ 是
Noether 环、$B$ 是 Noether $A$-代数，则映射 $A \to B$ 为正规态射，当且仅当
$B$ 是光滑 $A$-代数的直接极限。容易看出，Popescu 结果蕴含任意优良概形
上的 Artin 近似（优良性意味着对局部环 $A$，从 Hensel 化到完备化的映射
$A^{\text{h}} \to \hat A$ 为正规态射）。
论文给出 \cite{ArtinI} 主定理的一个适用于任意优良基概形、任意
$s \in S$ 的 ''groupoid'' 推广。因此 \cite{ArtinVersal} 中的结果在任意
优良基上成立。论文还讨论代数化的 \'etale-局部唯一性，以及对象的自同构群
是否自然作用在代数化的 Hensel 化上。
\end{quote}
\item Jason Starr: \emph{Artin's axioms, composition, and moduli spaces}
\cite{starr_artin}
\begin{quote}
论文证明 Artin 的代数化公理与 1-态射的复合相容。
\end{quote}
\item Martin Olsson: \emph{Deformation theory of representable
morphism of algebraic stacks} \cite{olsson_deformation}
\begin{quote}
本文用余切复形将概形态射的标准形变理论结果推广到代数叠的可表示态射。
这些结果不能视为 Illusie 一般理论的推论，因为可表示态射
$X \to \mathcal{X}$ 的余切复形并非由环化拓扑的态射定义（lisse-\'etale
site 不具有函子性）。
\end{quote}
\end{itemize}

\subsection{粗模空间}
\label{guide-subsection-coarse-moduli-spaces}

\noindent
讨论粗模空间的论文。
\begin{itemize}
\item Keel, Mori: \emph{Quotients in Groupoids} \cite{K-M}
\begin{quote}
长期以来，人们似乎以 ''folklore'' 形式知道分离 Deligne--Mumford 叠具有粗模空间。本文给出下述定理的严格（虽简略）证明：若 $\mathcal{X}$ 是 Noether 基概形上局部有限型的 Artin 叠，且惯性叠 $I_\mathcal{X} \to \mathcal{X}$ 有限，则存在粗模空间 $\phi : \mathcal{X} \to Y$，其中 $\phi$ 分离，$Y$ 是在 $S$ 上局部有限型的代数空间。惯性有限恰好是正确条件：存在分离的粗模空间 $\phi : \mathcal{X} \to Y$，其中 $\phi$ 分离，当且仅当惯性有限。
\end{quote}
\item Conrad: \emph{The Keel-Mori Theorem via Stacks} \cite{conrad}
\begin{quote}
Keel 与 Mori 的论文 \cite{K-M} 使用 groupoid 语言，读者往往觉得难以把握。Brian Conrad 给出了堆论版本的证明；虽然使用了精细的叠语言，却相当透明。Conrad 还去除了 Noether 假设。
\end{quote}
\item Rydh: \emph{Existence of quotients by finite groups and coarse moduli
spaces} \cite{rydh_quotients}
\begin{quote}
Rydh 去除了 \cite{K-M} 与 \cite{conrad} 中要求 $\mathcal{X}$ 在某个基上有限表示的假设。
\end{quote}
\item
Abramovich, Olsson, Vistoli: \emph{Tame stacks in positive characteristic}
\cite{tame}
\begin{quote}
作者将 \emph{tame Artin stack} 定义为具有有限惯性的 Artin 叠，并要求当
$\phi : \mathcal{X} \to Y$ 是粗模空间时，$\phi_*$ 在拟凝聚层上正合。他们证明，对有限惯性的 Artin 叠，下列条件等价：$\mathcal{X}$ tame；$\mathcal{X}$ 的稳定子群线性约化；以及在粗模空间上 \'etale 局部地，$\mathcal{X}$ 是仿射概形被线性约化群概形作用所得的商。对 tame Artin 叠，粗模空间具有特别良好的性质：例如粗模空间与任意基变换交换，而一般有限惯性 Artin 叠的粗模空间只与平坦基变换交换。
\end{quote}
\item Alper: \emph{Good moduli spaces for Artin stacks} \cite{alper_good}
\begin{quote}
对于具有无限仿射稳定子群（必然不分离）的一般 Artin 叠，粗模空间常常不存在。最简单的例子是
$[\mathbf{A}^1 / \mathbf{G}_m]$。本文将拟紧态射 $\phi : \mathcal{X} \to Y$ 称为
\emph{good moduli space}，若 $\mathcal{O}_Y \to \phi_* \mathcal{O}_\mathcal{X}$ 为同构，
且 $\phi_*$ 在拟凝聚层上正合。
这一概念既推广了 \cite{tame} 中的 tame Artin 叠，也概括了 Mumford 的几何不变理论：若 $G$ 是作用在线性空间 $X \subset \mathbf{P}^n$ 上的约化群，则半稳定轨迹的商叠到 GIT 商的态射 $[X^{ss}/G] \to X//G$ 是 good moduli space。该概念具有许多良好几何性质：
(1) $\phi$ 满射、普遍闭且普遍下沉；(2) $\phi$ 按闭包等价关系识别
$Y$ 与 $\mathcal{X}$ 中的点；(3) $\phi$ 对映到代数空间具有泛性质；
(4) good moduli space 在任意基变换下稳定；(5) Artin 叠上的向量丛下降到 good moduli space，当且仅当闭点处表示是平凡的。
\end{quote}
\end{itemize}


\subsection{交理论}
\label{guide-subsection-intersection-theory}

\noindent
讨论代数叠上交理论的论文。
\begin{itemize}
\item
Vistoli: \emph{Intersection theory on algebraic stacks and on their moduli
spaces} \cite{vistoli_intersection}
\begin{quote}
本文为 Deligne--Mumford 叠建立带有有理系数的交理论基础。若
$\mathcal{X}$ 是分离 Deligne--Mumford 叠，则带有有理系数的 Chow 群
$\CH_*(\mathcal{X})$ 定义为维数为 $k$ 的整闭子叠所生成的自由阿贝尔群模去有理等价。它具有平坦拉回、适当推前以及正则局部嵌入的广义 Gysin 同态。若 $\phi : \mathcal{X} \to Y$ 是模空间（即在几何点上为双射的适当态射），则存在诱导推前 $\CH_*(\mathcal{X}) \to \CH_k(Y)$，且为同构。
\end{quote}
\item Edidin, Graham: \emph{Equivariant Intersection Theory}
\cite{edidin-graham}
\begin{quote}
本文旨在为代数群 $G$ 作用于代数空间 $X$ 的商叠 $[X/G]$ 建立带整数系数的交理论，即建立 $X$ 的 $G$-等变交理论。仅用不变循环定义的等变 Chow 群不具备良好性质。作者推广 Totaro 对 $BG$ 的定义，并利用向量丛 $V \to X$ 给出自然同构 $\CH_i(X) \cong \CH_i(V)$，定义如下。记住定义中的 $\CH_i^G(X)$：
令 $\dim(X) = n$、$\dim(G) = g$。对每个 $i$，取一个 $l$ 维 $G$-表示 $V$，使 $G$ 在开集 $U \subset V$ 上自由作用，且补集的余维 $d > n - i$。于是 $X_G = [X \times U / G]$ 是代数空间（甚至可取为概形）。定义
$\CH_i^G(X) = \CH_{i + l - g}(X_G)$。对商叠定义
$\CH_i( [X/G]) = \CH_{i + g}^G(X) = \CH_{i + l}(X_G)$。
特别地，若 $i > \dim [X/G] = n - g$，则 $\CH_i([X/G]) = 0$，但当 $i \ll 0$ 时可能非零。例如
$\CH_i(B \mathbf{G}_m) = \mathbf{Z}$（$i \le 0$）。作者证明这些等变 Chow 群具有与普通 Chow 群相同的函子性；此外还证明若
$[X / G] \cong [Y / H]$，则 $\CH_i([X/G]) = \CH_i([Y/H])$，所以定义与叠如何表示为商叠无关。
\end{quote}
\item
Kresch: \emph{Cycle Groups for Artin Stacks} \cite{kresch_cycle}
\begin{quote}
Kresch 为任意 Artin 叠定义 Chow 群；在商叠情形，这一定义与 \cite{edidin-graham} 中 Edidin 和 Graham 的定义一致。对于具有仿射稳定子群的代数叠，该理论满足通常性质。
\end{quote}
\item Behrend and Fantechi: \emph{The intrinsic normal cone}
\cite{behrend-fantechi}
\begin{quote}
Behrend 与 Fantechi 将 Li 和 Tian 的构造推广到 Deligne--Mumford 叠，构造了虚基本类。
\end{quote}
\end{itemize}


\subsection{商叠}
\label{guide-subsection-quotient-stacks}

\noindent
商叠\footnote{在文献中，\emph{quotient stack} 常指形如 $[X/G]$ 的叠，其中 $X$ 为代数空间而 $G$ 为 $\text{GL}_n$ 的子群概形，而非任意平坦群概形。}
是 Artin 叠中非常重要的子类，涵盖代数几何学家研究的几乎所有模叠。商叠
$[X/G]$ 的几何就是 $X$ 的 $G$-等变几何。证明某性质对商叠成立通常更容易，
有些结果也只对商叠已知。下列论文讨论：代数叠何时是全局商叠？代数叠是否
在局部上是商叠？
\begin{itemize}
\item Laumon, Moret-Bailly: \cite[Chapter 6]{LM-B}
\begin{quote}
第 6 章包含关于代数叠局部与整体结构的若干事实。证明了：$S$ 上的代数叠
$\mathcal{X}$ 是商叠 $[Y/G]$（其中 $Y$ 可取代数空间、概形或仿射概形）
且 $G$ 为有限群，当且仅当存在代数空间、概形或仿射概形 $Y'$ 以及有限
\'etale 态射 $Y' \to \mathcal{X}$。还证明任意 $S$ 上的 Deligne--Mumford
叠及任意 $x : \Spec(K) \to \mathcal{X}$ 都有可表示、\'etale 且分离的态射
$\phi : [X/G] \to \mathcal{X}$，其中 $G$ 是作用于 $S$ 上仿射概形的有限群，
且满足 $\Spec(K) = [X/G] \times_\mathcal{X} \Spec(K)$。还详细讨论了具有几何连通纤维的展示的存在性。
\end{quote}
\item
Edidin, Hassett, Kresch, Vistoli:
\emph{Brauer Groups and Quotient stacks} \cite{ehkv}
\begin{quote}
首先，他们建立了关于代数叠何时为商叠的一些基础（虽不困难）事实（本文中
代数叠总假定在 Noether 概形上有限型）。对代数叠 $\mathcal{X}$，以下条件等价：
$\mathcal{X}$ 是商叠；存在向量丛 $V \to \mathcal{X}$，使每个几何点的稳定子群忠实作用于纤维；存在向量丛 $V \to \mathcal{X}$ 及局部闭子叠
$V^0 \subset V$，使 $V^0$ 可表示且满射到 $\mathcal{X}$。他们证明若存在来自代数空间的有限平坦覆盖，则代数叠是商叠；任意具有一般点平凡稳定子的光滑 Deligne--Mumford 叠都是商叠。
他们证明，Noether 概形 $X$ 上对应 $\beta \in H^2(X, \mathbf{G}_m)$ 的
$\mathbf{G}_m$-gerbe 是商叠，当且仅当 $\beta$ 在 Brauer 映射
$\text{Br}(X) \to \text{Br}'(X)$ 的像中。由此构造了一个不是商叠的非分离
Deligne--Mumford 叠。
\end{quote}
\item Totaro: \emph{The resolution property for schemes and stacks}
\cite{totaro_resolution}
\begin{quote}
若每个拟凝聚层都是某个向量丛的商，则叠具有 resolution property。第一主定理：
若 $\mathcal{X}$ 是闭点处稳定子群仿射的正规 Noether 代数叠，则下列等价：
(1) $\mathcal{X}$ 具有 resolution property；(2) $\mathcal{X} = [Y/\text{GL}_n]$，其中 $Y$ 拟仿射。
若 $\mathcal{X}$ 在域上有限型，则还等价于：(3)
$\mathcal{X} = [\Spec(A)/G]$，其中 $G$ 是在 $k$ 上有限型的仿射群概形。
商叠具有 resolution property 的蕴含由 Thomason 证明。第二主定理：
若 $\mathcal{X}$ 是域上的光滑 Deligne--Mumford 叠，惯性群
$I_\mathcal{X} \to \mathcal{X}$ 有限且一般点平凡，粗模空间是具有仿射对角的概形，则 $\mathcal{X}$ 具有 resolution property。另一个有趣结果是：若 $\mathcal{X}$ 是具有 resolution property 的 Noether 代数叠，则其对角仿射当且仅当闭点的稳定子群仿射。
\end{quote}
\item Kresch: \emph{On the Geometry of Deligne-Mumford Stacks}
\cite{kresch_geometry}
\begin{quote}
本文总结域上有限型 Deligne--Mumford 叠的一般结构结果，并包含关于商叠的有趣结果。证明任意光滑、分离、一般点 tame 且粗模空间拟射影的 Deligne--Mumford 叠都是商叠 $[Y/G]$，其中 $Y$ 拟射影、$G$ 为代数群。若 $\mathcal{X}$ 为粗模空间是概形的 Deligne--Mumford 叠，则 $\mathcal{X}$ 在 Zariski 局部为商叠，当且仅当它具有由概形被有限群商得到的叠的 Zariski 开覆盖。
若 $\mathcal{X}$ 是特征 0 域上 proper 的 Deligne--Mumford 叠、粗模空间为 $Y$，则：
$Y$ 射影且 $\mathcal{X}$ 为商叠，当且仅当 $Y$ 射影且 $\mathcal{X}$ 具有生成层，当且仅当 $\mathcal{X}$ 可闭嵌入一个具有射影粗模空间的光滑 proper DM 叠。由此启发定义：若 Deligne--Mumford 叠可闭嵌入具有射影粗模空间的光滑 proper Deligne--Mumford 叠，则称其为 \emph{projective}。
\end{quote}
\item Kresch, Vistoli \emph{On coverings of Deligne-Mumford stacks and
surjectivity of the Brauer map} \cite{kresch-vistoli}
\begin{quote}
在特征 0 且固定 $n$ 时，以下两条等价：(1) 每个 $n$ 维光滑 Deligne--Mumford
叠都是商叠；(2) $n$ 维光滑概形的 Azumaya Brauer 群等于上同调 Brauer 群。
\end{quote}
\item Kresch: \emph{Cycle Groups for Artin Stacks} \cite{kresch_cycle}
\begin{quote}
证明域上有限型且约化、稳定子群仿射的 Artin 叠具有按商叠给出的分层。
\end{quote}
\item Abramovich-Vistoli:
\emph{Compactifying the space of stable maps} \cite{abramovich-vistoli}
\begin{quote}
引理 2.2.3 证明：任意分离 Deligne--Mumford 叠在粗模空间上 \'etale 局部为
商叠 $[U/G]$，其中 $U$ 仿射、$G$ 为有限群。\cite[Theorem 2.12]{olsson_homstacks} 说明在此论证中 $G$ 甚至就是稳定子群。
\end{quote}
\item Abramovich, Olsson, Vistoli:
\emph{Tame stacks in positive characteristic} \cite{tame}
\begin{quote}
本文证明 tame Artin 叠在粗模空间上 \'etale 局部为仿射概形被稳定子群作用所得的商叠。
\end{quote}
\item Alper: \emph{On the local quotient structure of Artin stacks}
\cite{alper_quotient}
\begin{quote}
猜想：对 Artin 叠 $\mathcal{X}$ 与闭点 $x \in \mathcal{X}$，若稳定子群线性约化，
则存在带 $V$ 为代数空间的 \'etale 态射 $[V/G_x] \to \mathcal{X}$。文中给出
这一猜想的若干证据。基于 \cite{tame} 思想的简单形变理论论证表明，该结论在
形式局部成立。文章给出 Luna \'etale slice 定理的叠论证明：若
$\mathcal{X} = [\Spec(A)/G]$ 且 $G$ 线性约化，则在 GIT 商
$\Spec(A^G)$ 上 \'etale 局部地，$\mathcal{X}$ 是由稳定子群作商得到的商叠。
\end{quote}
\end{itemize}

\subsection{上同调}
\label{guide-subsection-cohomology}

\noindent
讨论代数叠上层的上同调的论文。
\begin{itemize}
\item Olsson: \emph{Sheaves on Artin stacks} \cite{olsson_sheaves}
\begin{quote}
本文发展拟凝聚层与可构造层理论，并证明基本上同调性质。它纠正了
\cite{LM-B} 关于 lisse-\'etale site 函子性的错误，并构造余切复形。此外证明：
proper 态射的 Grothendieck 基本定理、Grothendieck 存在定理、Zariski 连通性定理，
以及 proper 推前的凝聚层与可构造层有限性定理。
\end{quote}
\item Behrend: \emph{Derived $l$-adic categories for algebraic stacks}
\cite{behrend_derived}
\begin{quote}
证明代数叠的 Lefschetz 迹公式。
\end{quote}
\item Behrend: \emph{Cohomology of stacks} \cite{behrend_cohomology}
\begin{quote}
定义微分叠的 de Rham 上同调与拓扑叠的奇异上同调。
\end{quote}
\item Faltings:
\emph{Finiteness of coherent cohomology for proper fppf stacks}
\cite{faltings_finiteness}
\begin{quote}
证明 proper 态射下凝聚层直接像的凝聚性。
\end{quote}
\item Abramovich, Corti, Vistoli:
\emph{Twisted bundles and admissible covers} \cite{acv}
\begin{quote}
附录包含 tame Deligne--Mumford 叠的 \'etale 上同调 proper 基变换定理。
\end{quote}
\end{itemize}


\subsection{由概形给出的有限覆盖的存在性}
\label{guide-subsection-finite-covers}

\noindent
Deligne--Mumford 叠具有由概形给出的有限覆盖，是一项重要结果。在
Deligne--Mumford 叠的交理论中，它是为不可表示态射定义 proper 推前的关键。
关于 $\overline{\mathcal{M}}_g$ 的若干结果依赖于存在由\emph{光滑}概形给出的有限覆盖，
Looijenga 已证明了这一点。这个方向上较早的结果或许是
\cite[Theorem 6.1]{seshadri_quotients}，它处理了等变情形。
\begin{itemize}
\item Vistoli: \emph{Intersection theory on algebraic stacks and on their
moduli spaces} \cite{vistoli_intersection}
\begin{quote}
若 $\mathcal{X}$ 是具有模空间的 Deligne--Mumford 叠（即在几何点上为双射的
proper 态射），则存在来自概形 $X$ 的有限态射 $X \to \mathcal{X}$。
\end{quote}
\item Laumon, Moret-Bailly: \cite[Chapter 16]{LM-B}
\begin{quote}
Zariski 主定理的应用是：定理 16.6 证明，若 $\mathcal{X}$ 是 Noether 概形上
有限型的 Deligne--Mumford 叠，则存在有限、满射且一般为 \'etale 的态射
$Z \to \mathcal{X}$，其中 $Z$ 为概形。推论 16.6.2 还证明，任意 Noether
正规代数空间同构于有限群 $G$ 作用在正规概形 $X$ 上所得的代数空间商 $X'/G$。
\end{quote}
\item Edidin, Hassett, Kresch, Vistoli:
\emph{Brauer Groups and Quotient stacks} \cite{ehkv}
\begin{quote}
定理 2.7 表明：若 $\mathcal{X}$ 是 Noether 基概形 $S$ 上有限型的代数叠，
则对角 $\mathcal{X} \to \mathcal{X} \times_S \mathcal{X}$ 拟有限，当且仅当
存在来自概形 $X$ 的有限满射 $X \to F$。
\end{quote}
\item Kresch, Vistoli: \emph{On coverings of Deligne-Mumford stacks and
surjectivity of the Brauer map} \cite{kresch-vistoli}
\begin{quote}
证明：域上有限型、光滑且分离、粗模空间拟射影的 Deligne--Mumford 叠具有来自
光滑拟射影概形的有限平坦覆盖。
\end{quote}
\item Olsson: \emph{On proper coverings of Artin stacks} \cite{olsson_proper}
\begin{quote}
证明若 $\mathcal{X}$ 是 $S$ 上分离且有限型的 Artin 叠，则存在来自在 $S$ 上
拟射影的概形 $X$ 的 proper 满射 $X \to \mathcal{X}$。应用包括 proper 态射下
直接像层的凝聚性与可构造性，以及 proper Artin 叠的 Grothendieck 存在定理。
\end{quote}
\item Rydh: \emph{Noetherian approximation of algebraic spaces and stacks}
\cite{rydh_approx}
\begin{quote}
本文定理 B 如下：设 $X$ 是拟紧代数叠，具有拟有限分离对角（resp.\ 拟紧
Deligne--Mumford 叠具有拟紧分离对角）。则存在概形 $Z$ 及有限、有限表示且
满射的态射 $Z \to X$，它在稠密拟紧开子叠 $U \subset X$ 上是平坦的（resp.\ \'etale 的）。
\end{quote}
\end{itemize}


\subsection{刚化}
\label{guide-subsection-rigidification}

\noindent
刚化是从惯性中去除一个平坦子群的过程。例如，若 $X$ 是射影簇，
从 Picard 叠到 Picard 概形的态射就是对自同构群
$\mathbf{G}_m$ 的刚化。
\begin{itemize}
\item Abramovich, Corti, Vistoli:
\emph{Twisted bundles and admissible covers} \cite{acv}
\begin{quote}
设 $\mathcal{X}$ 是 $S$ 上的代数叠，$H$ 是 $S$ 上平坦、有限表示且分离的群概形。
假设对每个 $\xi \in \mathcal{X}(T)$ 都有嵌入
$H(T) \hookrightarrow \text{Aut}_{\mathcal{X}(T)}(\xi)$，并且在拉回下相容：
对每个覆盖 $f: T \rightarrow T'$ 上的箭头 $\phi : \xi \rightarrow \xi'$ 与
$g \in H(T')$，都有 $g \circ \phi = \phi \circ f^*g$。
则存在代数叠 $\mathcal{X}/H$ 及态射
$\rho : \mathcal{X} \rightarrow \mathcal{X}/H$，它是 fppf gerbe，并且对每个
$\xi \in \mathcal{X}(T)$，态射
$\text{Aut}_{\mathcal{X}(T)} (\xi) \rightarrow
\text{Aut}_{\mathcal{X}/H (T)} (\xi)$
为满射，核为 $H(T)$。
\end{quote}
\item Romagny: \emph{Group actions on stacks and applications}
\cite{romagny_actions}
\begin{quote}
讨论群作用相对于刚化的行为。
\end{quote}
\item Abramovich, Graber, Vistoli:
\emph{Gromov-Witten theory for Deligne-Mumford stacks} \cite{agv}
\begin{quote}
附录按照 \cite{acv} 总结刚化，并给出两种替代解释。本文还包含沿闭子叠粘合代数叠以及对线丛取根的构造。
\end{quote}
\item
Abramovich, Olsson, Vistoli: \emph{Tame stacks in positive characteristic}
(\cite{tame})
\begin{quote}
附录处理更复杂的情形：惯性中的平坦子叠
$H \subset I_\mathcal{X}$ 是正规但不一定中心的。
\end{quote}
\end{itemize}

\subsection{叠曲线}
\label{guide-subsection-stacky-curves}

\noindent
讨论叠曲线的论文。
\begin{itemize}
\item Abramovich, Vistoli: \emph{Compactifying the space of stable maps}
\cite{abramovich-vistoli}
\begin{quote}
本文引入 \emph{twisted curves}。从稳定曲线到代数叠的稳定映射的模空间通常并不紧。
作者利用扭曲曲线构造模叠；当目标是 tame Deligne--Mumford 叠且粗模空间射影时，
该模叠是 proper 的。
\end{quote}
\item Behrend, Noohi: \emph{Uniformization of Deligne-Mumford curves}
\cite{behrend-noohi}
\begin{quote}
证明 Deligne--Mumford 解析曲线的统一化定理。
\end{quote}
\end{itemize}

\subsection{Hilbert、Quot、Hom 与分支簇叠}
\label{guide-subsection-hilbert-quot-hom}

\noindent
讨论 Hilbert 概形及类似对象的论文。
\begin{itemize}
\item Vistoli: \emph{The Hilbert stack and the theory of moduli of families}
\cite{vistoli_hilbert}
\begin{quote}
若 $\mathcal{X}$ 是在局部 Noether 且局部分离代数空间 $S$ 上分离、局部有限型的代数叠，
Vistoli 定义 Hilbert 叠 $\mathcal{H}\text{ilb}(\mathcal{F} / S)$，参数化来自 proper 概形的有限且非分歧态射。论文未证明地声称 $\mathcal{H}\text{ilb}(\mathcal{F} / S)$ 是代数叠。由此证明：若 $\mathcal{X}$ 如上，则当 $T$ 在 $S$ 上 proper 且平坦时，Hom 叠 $\mathcal{H} \text{om}_S(T, \mathcal{X})$ 是代数叠。
\end{quote}
\item Olsson, Starr: \emph{Quot functors for Deligne-Mumford stacks}
\cite{olsson-starr}
\begin{quote}
若 $\mathcal{X}$ 是代数空间 $S$ 上分离且局部有限表示的 Deligne--Mumford 叠，
$\mathcal{F}$ 是局部有限表示的 $\mathcal{O}_\mathcal{X}$-模，则 Quot 函子
$\text{Quot}(\mathcal{F} / \mathcal{X} / S)$ 由一个在 $S$ 上分离且局部有限表示的代数空间表示。
本文还定义生成层，并证明对由有限群作用于概形得到的全局商 tame、分离
Deligne--Mumford 叠存在生成层。
\end{quote}
\item Olsson: \emph{Hom-stacks and Restrictions of Scalars}
\cite{olsson_homstacks}
\begin{quote}
设 $\mathcal{X}$ 与 $\mathcal{Y}$ 是代数空间 $S$ 上局部有限表示、对角有限的 Artin 叠，
其中 $\mathcal{X}$ 在 $S$ 上 proper 且平坦，并且在 $S$ 上 fppf 局部地，
$\mathcal{X}$ 有来自代数空间的有限有限表示平坦覆盖（例如 $\mathcal{X}$ 是
Deligne--Mumford 叠或 tame Artin 叠）。则 $\Hom_S(\mathcal{X}, \mathcal{Y})$
是 $S$ 上局部有限表示的 Artin 叠。
\end{quote}
\item Alexeev and Knutson: \emph{Complete moduli spaces of branchvarieties}
(\cite{alexeev-knutson})
\begin{quote}
他们将 $\mathbf{P}^n$ 的 branchvariety 定义为来自\emph{约化}概形 $X$ 的有限态射
$X \rightarrow \mathbf{P}^n$。证明固定 Hilbert 多项式及 $i$ 维分支的总次数后，
branchvariety 的模叠是具有有限稳定子的 proper Artin 叠。并比较 branchvariety 叠
与 Hilbert 概形、Chow 概形和稳定映射模空间。
\end{quote}
\item Lieblich: \emph{Remarks on the stack of coherent algebras}
\cite{lieblich_remarks}
\begin{quote}
本文通过把叠构造成概形 $Y$ 的结构层上的代数叠，构造 Alexeev--Knutson
branch-variety 叠在概形 $Y$ 上的推广，并给出 $\text{Quot}$ 与 $\Hom$ 空间的存在性证明。
\end{quote}
\item Starr: \emph{Artin's axioms, composition, and moduli spaces}
\cite{starr_artin}
\begin{quote}
作为主结果的应用，给出 Vistoli Hilbert 叠 \cite{vistoli_hilbert} 与 Alexeev--Knutson
branchvariety 叠 \cite{alexeev-knutson} 的共同推广。若 $\mathcal{X}$ 是优良概形
$S$ 上局部有限型、对角有限的代数叠，则参数化从带有 $G$-ample 线丛 $L$ 的
proper 代数空间 $T$ 到 $\mathcal{X}$ 的态射 $g: T \rightarrow \mathcal{X}$ 的叠
$\mathcal{H}$ 是 $S$ 上局部有限型 Artin 叠。
\end{quote}
\item Lundkvist and Skjelnes:
\emph{Non-effective deformations of Grothendieck's Hilbert functor}
\cite{lundkvist-skjelnes}
\begin{quote}
证明非分离概形的 Hilbert 函子不可表示，因为存在非有效形变。
\end{quote}
\item Halpern-Leistner and Preygel:
\emph{Mapping stacks and categorical notions of properness}
\cite{HL-P}
\begin{quote}
本文证明，在比 Olsson 论文更一般或至少不同的源与目标假设下，Hom 叠仍是代数的。
\end{quote}
\end{itemize}



\subsection{环面叠}
\label{guide-subsection-toric}

\noindent
环面叠提供了大量例子，也是检验猜想的自然场所：环面叠的几何与其叠扇的组合学
之间存在字典，正如环面簇在概形理论中提供例子与反例一样。
\begin{itemize}
\item Borisov, Chen and Smith: \emph{The orbifold Chow ring of toric
Deligne-Mumford stacks} \cite{bcs}
\begin{quote}
受 Cox 对环面簇构造的启发，本文将光滑环面 DM 叠定义为与称为
\emph{stacky fan} 的组合对象相联系的显式商叠。
\end{quote}
\item Iwanari: \emph{The category of toric stacks} \cite{iwanari_toric}
\begin{quote}
本文将 \emph{toric triple} 定义为光滑 Deligne--Mumford 叠
$\mathcal{X}$，带有稠密像的开浸入 $\mathbf{G}_m \hookrightarrow \mathcal{X}$
（因此 $\mathcal{X}$ 是 orbifold）以及作用
$\mathcal{X} \times \mathbf{G}_m \rightarrow \mathcal{X}$。证明 toric triple
的 2-范畴与 stacky fan 的 1-范畴等价，并讨论 toric triple 与
\cite{bcs} 中光滑环面 DM 叠定义的关系。
\end{quote}
\item Iwanari: \emph{Integral Chow rings for toric stacks}
\cite{iwanari_chow}
\begin{quote}
将 Cox 对环面簇的 $\Delta$-collections 推广到环面 orbifold。
\end{quote}
\item Perroni: \emph{A note on toric Deligne-Mumford stacks}
\cite{perroni}
\begin{quote}
将 Cox 的 $\Delta$-collections 及 Iwanari 的论文
\cite{iwanari_chow} 推广到一般光滑环面 DM 叠。
\end{quote}
\item Fantechi, Mann, and Nironi: \emph{Smooth toric DM stacks}
\cite{fmn}
\begin{quote}
本文将光滑环面 DM 叠定义为带有 DM 环面 $\mathcal{T}$ 作用的光滑 DM 叠
$\mathcal{X}$（即 Picard 叠同构于 $T \times BG$，其中 $G$ 有限），且具有同构于
$\mathcal{T}$ 的开稠密轨道。作者给出“自底向上”的描述，并证明光滑环面 DM 叠与
stacky fan 等价。
\end{quote}
\item Geraschenko and Satriano: \emph{Toric Stacks I and II}
\cite{gs_toric1} and \cite{gs_toric2}
\begin{quote}
这两篇论文将环面叠定义为环面簇被其环面的某个子群作商得到的叠。generically
stacky toric stack 定义为环面叠中的环面不变子叠；该定义包含并推广了早期定义。
第一篇建立叠扇组合学与相应叠几何之间的字典，并给出光滑环面叠的模解释，
推广 \cite{perroni} 中的解释。第二篇证明环面叠的内在刻画。
\end{quote}
\end{itemize}



\subsection{形式函数定理与 Grothendieck 存在定理}
\label{guide-subsection-theorem-formal-functions}

\noindent
这些论文推广形式函数定理
\cite[III.4.1.5]{EGA}（有时称为 proper 态射的 Grothendieck 基本定理）
以及 Grothendieck 存在定理 \cite[III.5.1.4]{EGA}。
\begin{itemize}
\item Knutson: \emph{Algebraic spaces} \cite[Chapter V]{Kn}
\begin{quote}
将这些定理推广到代数空间。
\end{quote}
\item Abramovich-Vistoli: \emph{Compactifying the space of stable maps}
\cite[A.1.1]{abramovich-vistoli}
\begin{quote}
将这些定理推广到 tame Deligne--Mumford 叠。
\end{quote}
\item Olsson and Starr: \emph{Quot functors for Deligne-Mumford stacks}
\cite{olsson-starr}
\begin{quote}
将这些定理推广到分离 Deligne--Mumford 叠。
\end{quote}
\item Olsson: \emph{On proper coverings of Artin stacks}
\cite{olsson_proper}
\begin{quote}
给出到 proper Artin 叠的推广。
\end{quote}
\item Conrad: \emph{Formal GAGA on Artin stacks} \cite{conrad_gaga}
\begin{quote}
给出到 proper Artin 叠的推广，并证明形式 GAGA 定理。
\end{quote}
\item Olsson: \emph{Sheaves on Artin stacks} \cite{olsson_sheaves}
\begin{quote}
给出到 proper Artin 叠的推广的另一种证明。
\end{quote}
\end{itemize}


\subsection{叠上的群作用}
\label{guide-subsection-group-actions}

\noindent
代数叠上的群作用自然出现。例如，对称群 $S_n$ 作用于
$\overline{\mathcal{M}}_{g, n}$；若群 $G$ 作用于概形 $X$，则
$\text{Aut}(X)$ 中 $G$ 的正规化子作用于 $[X/G]$。此外，环面作用经常出现在
Gromov--Witten 理论中。
\begin{itemize}
\item Romagny: \emph{Group actions on stacks and applications}
\cite{romagny_actions}
\begin{quote}
本文精确说明群如何作用于代数叠，并证明固定点的存在性以及群概形作用在代数叠上时商的存在性。另见 Romagny 较早的笔记 \cite{romagny_notes}。
\end{quote}
\end{itemize}

\subsection{对线丛取根}
\label{guide-subsection-root-stacks}

\noindent
这一有用构造由 Cadman 以及 Abramovich、Graber、Vistoli 独立发现。给定带有效 Cartier 除子
$D$ 的概形 $X$，$r$ 次根叠是沿 $D$ 在 $X$ 上分歧的 Artin 叠，在 $D$ 上有
$\mu_r$ 稳定子群，而远离 $D$ 时类似概形。
\begin{itemize}
\item Charles Cadman
\emph{Using Stacks to Impose Tangency Conditions on Curves}
\cite{cadman}
\item Abramovich, Graber, Vistoli: \emph{Gromov-Witten theory for
Deligne-Mumford stacks} \cite{agv}
\end{itemize}

\subsection{其他论文}
\label{guide-subsection-other}

\noindent
其他论文的杂集。
\begin{itemize}
\item Lieblich: \emph{Moduli of twisted sheaves} \cite{lieblich_twisted}
\begin{quote}
本文概述 gerbe 与扭曲层。若 $\mathcal{X} \rightarrow X$ 是 $\mu_n$-gerbe，
$X$ 是具有光滑连通几何纤维的射影相对曲面，则证明 $\mathcal{X}$-扭曲半稳定层的
叠是在 $S$ 上局部有限表示的 Artin 叠。本文还发展 Artin 叠上伴随点与层纯性的理论。
\end{quote}
\item Lieblich, Osserman:
\emph{Functorial reconstruction theorem for stacks}
\cite{lieblich-osserman}
\begin{quote}
证明代数叠何时可以由其伴随函子重构的一些出人意料且有趣的结果。
\end{quote}
\item David Rydh:
\emph{Noetherian approximation of algebraic spaces and stacks}
\cite{rydh_approx}
\begin{quote}
证明每个具有拟有限对角的拟紧代数叠都可由 Noether 叠近似。应用包括去除
Chevalley、Serre、Zariski 与 Chow 结果中的 Noether 假设。
\end{quote}
\end{itemize}



\section{其他领域中的叠}
\label{guide-section-stacks-areas}

\begin{itemize}
\item Behrend and Noohi: \emph{Uniformization of Deligne-Mumford curves}
\cite{behrend-noohi}
\begin{quote}
概述拓扑叠、解析叠与代数叠，并进行比较。
\end{quote}
\item Behrang Noohi: \emph{Foundations of topological stacks I} \cite{noohi}
\item David Metzler: \emph{Topological and smooth stacks} \cite{metzler}
\end{itemize}


\section{高阶叠}
\label{guide-section-higher-stacks}

\begin{itemize}
\item Lurie: \emph{Higher topos theory} \cite{lurie_topos}
\item Lurie: \emph{Derived Algebraic Geometry I - V}
\cite{dag1}, \cite{dag2}, \cite{dag3}, \cite{dag4}, \cite{dag5}
\item To\"en: \emph{Higher and derived stacks: a global overview}
\cite{toen_higher}
\item To\"en and Vezzosi: \emph{Homotopical algebraic geometry I, II}
\cite{hag1}, \cite{hag2}
\end{itemize}
